\documentclass[11pt,a4paper]{article} % Remove 'draft' once all images are uploaded
\usepackage[utf8]{inputenc}
\usepackage{amsmath, amssymb, amsthm}
\usepackage{graphicx}
\usepackage{hyperref}
\usepackage{geometry}
\usepackage{booktabs}
\usepackage{caption}
\usepackage{subcaption}
\sloppy 

\geometry{margin=1in}

\title{\textbf{SAiDL Summer Assignment 2026: Hardening Transformer-Based Reinforcement Learning for Partial Observability}}
\author{Gautham [REDACTED]}
\date{\today}

\begin{document}

\maketitle

\begin{abstract}
This report details the implementation and optimization of a Causal Transformer-based Reinforcement Learning agent for the MuJoCo Hopper-v5 environment. We investigate the agent's ability to overcome Partial Observability (POMDP) through temporal attention mechanisms. Our results demonstrate that the Transformer architecture significantly outperforms standard MLP-based TD3 agents in environments with sensory failure, observational noise, and credit assignment delays. We further refine the agent through a Reinforcement Learning from Human Feedback (RLHF) pipeline utilizing a Jury Ensemble of reward models.
\end{abstract}

\section{Introduction}
Reinforcement Learning (RL) agents often fail in real-world scenarios due to the non-Markovian nature of environments (POMDPs). Standard architectures like Multi-Layer Perceptrons (MLPs) lack the memory required to reconstruct state variables from historical observations. This project explores the use of a Causal Transformer to solve these challenges.

\section{Task 1: Transformer Architectures for Locomotion}
\subsection{The Foundation: Twin Delayed DDPG (TD3)}
Before implementing the Transformer, we established a robust Reinforcement Learning framework using \textbf{Twin Delayed DDPG (TD3)}. TD3 improves upon standard DDPG by addressing the overestimation bias of Q-values through clipped double-Q learning, target policy smoothing, and delayed policy updates.

\subsection{Initial Challenges: The Unoptimized Baseline}
Our initial experiments with the Transformer architecture exhibited significant instability and slow convergence. As shown in Figure \ref{fig:before_comparison}, the early iterations suffered from high variance and a low reward ceiling, often failing to maintain a stable gait beyond 200 points.

\begin{figure}[!ht]
    \centering
    \begin{subfigure}[b]{0.48\textwidth}
        \centering
        \includegraphics[width=\textwidth]{Old Graphs/TD3_Baseline_Results_Old.png}
        \caption{Initial MLP Baseline (Seed 0-2).}
    \end{subfigure}
    \hfill
    \begin{subfigure}[b]{0.48\textwidth}
        \centering
        \includegraphics[width=\textwidth]{Old Graphs/Transformer_Comparison_Old.png}
        \caption{Unoptimized Transformer Ablation.}
    \end{subfigure}
    \caption{Phase 1 (The Before): Results from the early stages. Note the erratic convergence.}
    \label{fig:before_comparison}
\end{figure}

\subsection{Engineering Optimizations and Stability Fixes}
To overcome these hurdles, we implemented a series of industrial-grade optimizations:
\begin{itemize}
    \item \textbf{AMP (Automatic Mixed Precision)}: We integrated \texttt{GradScaler} to allow for 16-bit precision training, accelerating attention head learning.
    \item \textbf{Global Normalization Strategy}: We shifted to a global \textbf{RunningMeanStd} algorithm to stabilize the Transformer's input distribution.
    \item \textbf{Causal History Buffer}: We re-engineered the experience replay logic to treat the environment as a sequence of transitions, preventing future leakage.
\end{itemize}

\subsection{Final Results: The Optimized Plateau}
As shown in Figure \ref{fig:task1_comparison}, these optimizations resulted in a dramatic increase in performance. The agent now achieves a stable plateau of $\sim$600 points.

\begin{figure}[!ht]
    \centering
    \begin{subfigure}[b]{0.48\textwidth}
        \centering
        \includegraphics[width=\textwidth]{Graphs/TD3_Baseline_Results.png}
        \caption{Final MLP Baseline.}
        \label{fig:mlp_baseline}
    \end{subfigure}
    \hfill
    \begin{subfigure}[b]{0.48\textwidth}
        \centering
        \includegraphics[width=\textwidth]{Graphs/Transformer_Comparison.png}
        \caption{Optimized Transformer Ablation.}
        \label{fig:transformer_comparison}
    \end{subfigure}
    \caption{Phase 2 (The After): Final performance comparison after implementing optimizations.}
    \label{fig:task1_comparison}
\end{figure}

\section{Task 2: The Robustness Frontier (POMDP)}
\subsection{2a \& 2b: Sensor Failure and Noise (The Blindfold Test)}
In the first stressor, we disabled all velocity-related sensors. Observations were masked to provide only positional and joint angle data, with added Gaussian noise ($\sigma=0.1$).

\subsection{2c: Credit Assignment Delay (The Lag Test)}
The second stressor introduced a 10-frame observation lag, creating a "Temporal Disconnect" between actions and sensory responses.

\begin{figure}[!ht]
    \centering
    \includegraphics[width=0.7\textwidth]{Graphs/Combined_POMDP_Final.png}
    \caption{The Triple-Threat Victory: Transformer performance under simultaneous stressors.}
    \label{fig:combined_pomdp}
\end{figure}

\subsection{2d: Reinforcement Learning from Human Feedback (RLHF)}
We implemented an RLHF pipeline utilizing a \textbf{Jury of three independent Reward Models}.
\begin{equation}
R_{Jury}(s, a) = \min(RM_1, RM_2, RM_3)
\end{equation}

\begin{figure}[!ht]
    \centering
    \includegraphics[width=0.7\textwidth]{Graphs/RLHF_Final_Comparison.png}
    \caption{RLHF Alignment: Performance against the Optimized MLP baseline.}
    \label{fig:rlhf_final}
\end{figure}

\section{Task 3: Diagnostic Analysis and Positional Encodings}
\subsection{Attention Saliency and Attribution}
We implemented a \textbf{Saliency Attribution Engine} and found that the agent exhibits "Action Shots"---sharp focus on the moments exactly 0.5s prior.

\begin{figure}[!ht]
    \centering
    \includegraphics[width=0.7\textwidth]{Graphs/Attribution_Comparison.png}
    \caption{Saliency Attribution: Spikes indicate critical historical frames.}
    \label{fig:attribution}
\end{figure}

\subsection{The Temporal Ruler: RoPE vs. Sinusoidal Encodings}
We benchmarked Rotary Positional Embeddings (RoPE) against Sinusoidal Encodings. Sinusoidal Encodings provided the most stable "Temporal Ruler."

\begin{figure}[!ht]
    \centering
    \includegraphics[width=0.7\textwidth]{Graphs/PosEnc_Ablation.png}
    \caption{Positional Encoding Impact: Sinusoidal encodings yielded higher stability.}
    \label{fig:pos_enc_ablation}
\end{figure}

\subsection{Conclusion: The Surgical Calculus}
Our research demonstrates that a Causal Transformer equipped with Sinusoidal Positional Encodings is capable of performing "Surgical Calculus" on its own history, mathematically deriving required derivatives for balance.

\end{document}
